\documentclass[10pt,a4paper]{amsart}
\usepackage[margin=19mm]{geometry}
\usepackage{amsmath,amssymb,mathtools}
\usepackage[scr=rsfs,cal=euler]{mathalfa}
\usepackage{xfrac}
\usepackage[unicode,hidelinks,bookmarks=false]{hyperref}
\newcommand{\ie}{i.e.\ }
\newcommand{\cf}{cf.\ }
\newcommand{\resp}{respectively\ }
\newcommand{\Subsection}{\subsection}
\providecommand{\nobreakdash}{\nobreak\hbox{-}}
\setlength{\emergencystretch}{2em}
\multlinegap0pt
\newtheorem{lemma}{Lemma}
\newcommand{\K}{\mathcal{K}}
\renewcommand{\L}{\mathcal{L}}
\newcommand{\Z}{\mathbb{Z}}
\providecommand{\ba}{\bar a}
\providecommand{\bb}{\bar b}
\providecommand{\tightlist}{%
  \setlength{\itemsep}{0pt}\setlength{\parskip}{0pt}}
\title{Scott sentence complexities of linear orderings}
\author{David Gonzalez, Dino Rossegger}
\date{Source: arXiv:2305.07126v1 (CC BY 4.0)}
\begin{document}
\maketitle

\noindent\emph{Source and licence.} Scott sentence complexities of linear orderings. Version arXiv:2305.07126v1 (CC BY 4.0), distributed by arXiv under the Creative Commons Attribution 4.0 International licence, http://creativecommons.org/licenses/by/4.0/, as recorded in the arXiv licence metadata for this exact version and shown on the abstract page https://arxiv.org/abs/2305.07126v1. Full licence notice: \texttt{licenses/arxiv-2305.07126-CC-BY-4.0.txt}.\par\smallskip

\noindent\textit{Editorial scope.} Counted unit: the article's lemma lem:timesZ (source0.tex lines 1147--1158). The article's complete original proof of it is delivered with it (source0.tex lines 1160--1266, one \texttt{proof} environment of 107 lines). No mathematics is written, repaired or re-derived: the statement and the proof are the article's own lines, in the article's own notation, and no retained line is substituted or re-flowed.  No further source-local context is needed: every cross-reference of every retained line resolves inside the two delivered spans.  Nothing was omitted from any retained span, so the package carries no omission record.  No retained line contains a citation, so the wrapper carries no bibliography and no bibliography entry.  No retained line contains a figure environment, an image-inclusion command or drawing code, so no image is delivered and the package declares no asset dependency.  Editorial formatting only, in addition: an AMS \texttt{amsart} preamble that restates the article's own declarations and macros used by the retained lines, namely the environments \texttt{lemma} on separate counters, together with the standard packages those lines need.  The retained environments print in this document as Lemma 1 in order of appearance, whereas the article prints the same items with the numbering of its own full document.  The complete native source of the article as distributed in the arXiv e-print for this exact version is bundled unchanged as \texttt{sources/141394/source0.tex}.
\begin{lemma}\label{lem:timesZ}
For the sake of organization, consider the following ordinal indexed propositions.
\begin{enumerate}\tightlist
\item[($A_\alpha$)] For any $\K$ and $\L$, $\Z^\alpha\cdot \K
    \leq_{2\alpha}\Z^\alpha\cdot \L$.
\item[($B_\alpha$)] For any $\K$ and $\L$ with $|\K|\geq|L|$, $\Z^\alpha\cdot \K
    \leq_{2\alpha+1}\Z^\alpha\cdot \L$.
\item[($C_\alpha$)] For any $\K$ and any $\L$ without a last element,
    $\Z^\alpha\cdot(\omega+\K)\leq_{2\alpha+2} \Z^\alpha\cdot(\omega+\L)$.
\end{enumerate}
For any countable ordinal $\alpha$, $A_\alpha$, $B_\alpha$ and $C_\alpha$ are true.
\end{lemma}

\begin{proof}
    The proof is by induction with the statements $A_0$ and $B_0$ trivially true.
    We start by showing that $B_\alpha$ implies $A_{\alpha+1}$. 

    Select a tuple $a_0\cdots a_n$ in $\Z^{\alpha+1}\cdot \L=\Z^\alpha\cdot\Z\cdot
    \L$. We will find a tuple $b_0\dots b_n$ in $\Z^{\alpha+1}\cdot
    \K=\Z^\alpha\cdot\Z\cdot \K$ such that $(\Z^{\alpha+1}\cdot \L,\ba
    )\leq_{2\alpha+1} (\Z^{\alpha+1}\cdot \K,\bb)$. Pick a point $b_0$ arbitrarily.
    Given $b_j$, we pick $b_{j+1}$ as follows. If $a_j$ and $a_{j+1}$ are in the
    same copy of $\Z^{\alpha+1}$, we can pick $b_{j+1}$ such that
    $(a_j,a_{j+1})\cong(b_j,b_{j+1})$. If not, we pick $b_{j+1}$
    arbitrarily in the copy of $\Z^{\alpha}$ that is the successor of the successor
    of the copy of $b_j$. In other words, we pick $b_{j+1}$ with  $(b_j,b_{j+1})\cong \zeta^*_\alpha +\Z^\alpha+\zeta_\alpha$.
    
To show that our choice of $\bb$ is correct, we have to confirm the following
relations between non-isomorphic intervals.
\begin{enumerate}\tightlist
\item $(b_j,b_{j+1})\cong \zeta^*_\alpha +\Z^\alpha+\zeta_\alpha \geq_{2\alpha+1} \zeta^*_\alpha
    +\Z^\alpha\cdot \L_1+\zeta_\alpha$, where $\L_1$ is an infinite order,
\item $(-\infty,b_0)\cong \Z^\alpha\cdot \K_1+\zeta_\alpha \geq_{2\alpha+1}\Z^\alpha\cdot
    \L_1+\zeta_\alpha$, where $\L_1$ and $\K_1$ are some initial segments of the $\Z^\alpha$, and therefore infinite, or
\item $(b_{n},\infty)\cong\zeta^*_\alpha +\Z^\alpha\cdot \L_1 \geq_{2\alpha+1} \zeta^*_\alpha
    +\Z^\alpha\cdot \K_1$, where $\L_1$ and $\K_1$ are some final segments of the $\Z^\alpha$, and therefore infinite.
\end{enumerate} 
It follows immediately from $B_\alpha$ that all of these back-and-forth relations hold.

We now show that $B_\alpha$ implies $C_\alpha$. So assume that $\K$ and $\L$ have
no last element and that $B_\alpha$ holds. 
Select a tuple $a_0\cdots  a_n$ in $\Z^\alpha\cdot(\omega+\L)$.
Without loss of generality, $a_0$ is in the first copy of $\Z^\alpha$. We will
find a tuple $b_0\dots b_n$ in $\Z^\alpha\cdot(\omega+\K)$ such that
$(\Z^\alpha\cdot (\omega+\L),\ba)\leq_{2\alpha+1}
(\Z^\alpha\cdot(\omega+\K),\bb)$. For the $a_j$ in the initial $\Z^\alpha\cdot\omega$, the $b_j$ are picked isomorphically. 
We now define the rest of the $b_j$ by induction so that all of them are in the initial $\Z^\alpha\cdot\omega$.
If $a_j$ is in the same $\Z^\alpha$ block as $a_{j+1}$ or $a_{j+1}$ is in the
succesor of $a_j$'s block, define $b_{j+1}$ so that
$(b_j,b_{j+1})\cong (a_j,a_{j+1})$.
Otherwise, define $b_{j+1}$ as being some element in the successor of the
successor of the $\Z^\alpha$ block of $b_j$. 
Not including isomorphic intervals, we only need to check the following cases. 
\begin{enumerate}
    \item $(b_j,b_{j+1})\cong\zeta^*_\alpha +\Z^\alpha+\zeta_\alpha \geq_{2\alpha+1} \zeta^*_\alpha
    +\Z^\alpha\cdot \L_1+\zeta_\alpha$ where $\L_1$ is infinite, 
\item $(b_n,\infty)\cong\zeta^*_\alpha +\Z^\alpha\cdot (\omega+\K) \geq_{2\alpha+1}
    \zeta^*_\alpha +\Z^\alpha\cdot \L_1$, where $\L_1$ is some final segment of
    $\omega+\L$, and thus infinite.
\end{enumerate} 
All of these relations follow from $B_\alpha$, as desired.

To complete the proof of the successor step of the induction, we show that
$C_\alpha$ and $A_{\alpha+1}$ imply $B_{\alpha+1}$.

Select a tuple  $a_0\cdots a_n$ in $\Z^{\alpha+1}\cdot \L$. These points belong
to at most $n$ distinct copies of $\Z^{\alpha+1}$. Because $|\K|\geq| \L|$, we
can match each chosen copy of $\Z^{\alpha+1}$ to one in $\Z^{\alpha+1}\cdot \K$.
Furthermore, if there is a copy of $\Z^{\alpha+1}$ between two of the chosen
copies in $\Z^{\alpha+1}\cdot \L$, we can assure that the corresponding copies
in $\Z^{\alpha+1}\cdot \K$ are also not successors. Also, if $a_n$ is not in the
last copy of $\Z^{\alpha+1}$, we can assure that neither is $b_n$ and, similarly
for $a_0$ and $b_0$ in the first copy. In particular, not including the
isomorphic intervals, we only need to check the following cases.
\begin{enumerate}
    \item $(b_j,b_{j+1})\cong \zeta^*_{\alpha+1} +\Z^{\alpha+1}\cdot
        \K_1+\zeta_{\alpha+1}\geq_{2\alpha+2}\zeta^*_{\alpha+1}
        +\Z^{\alpha+1}\cdot \L_1+\zeta_{\alpha+1}$ for some orders $\K_1$ and $\L_1$,
    \item $(b_n,\infty)\cong\zeta^*_{\alpha+1} +\Z^{\alpha+1}\cdot
        \K_1\geq_{2\alpha+2}\zeta^*_{\alpha+1} +\Z^{\alpha+1}\cdot \L_1$ for
        some orders $\K_1$ and $\L_1$,
\item $(-\infty,b_0)\cong\Z^{\alpha+1}\cdot
    \K_1+\zeta_{\alpha+1}\geq_{2\alpha+2}\Z^{\alpha+1}\cdot
    \L_1+\zeta_{\alpha+1}$ for some orders $\K_1$ and $\L_1$,
\item $(-\infty,b_0)\cong\zeta_{\alpha+1}\geq_{2\alpha+2} \Z^{\alpha+1}\cdot
    \L_1+\zeta_{\alpha+1}$ for some order $\L_1$,
\item $(b_n,\infty)\cong\zeta^*_{\alpha+1} +\Z^{\alpha+1}\cdot
    \K_1\geq_{2\alpha+2} \zeta^*_{\alpha+1} $ for some order $\K_1$.
\end{enumerate}

The first three cases are handled immediately by $A_{\alpha+1}$. By symmetry it
is sufficient to show case (5). However, this is the same as showing that 
$$\zeta_{\alpha}^* + \Z^\alpha\cdot\omega \leq_{2\alpha+2}\zeta_{\alpha}^* +
\Z^{\alpha}\cdot (\omega+\Z\cdot \K_1).$$
This follows directly from $C_\alpha$.




At last we consider the remaining limit levels. We already have seen that
$B_\lambda$ implies $C_\lambda$. The statement $A_\lambda$ follows immediately from $A_\gamma$ for $\gamma<\lambda$. 
A bit more nuanced is the issue of $B_\lambda$. It follows from $A_\lambda$ and $C_\gamma$ for $\gamma<\lambda$.
Analogous to the proof in the successor case, we analyze the possibilities:
\begin{enumerate}
\item $\zeta^*_{\lambda} +\Z^{\lambda}\cdot
    \K_1+\zeta_{\lambda}\geq_\lambda \zeta^*_{\lambda} +\Z^{\lambda}\cdot \L_1+\zeta_{\lambda} $ for some orders $\L_1$ and $\K_1$,
\item $\zeta^*_{\lambda} +\Z^{\lambda}\cdot \K_1\geq_\lambda \zeta^*_{\lambda} +\Z^{\lambda}\cdot \L_1 $ for some orders $\L_1$ and $\K_1$,
\item $\Z^{\lambda}\cdot \K_1+\zeta_{\lambda}\geq_\lambda \Z^{\lambda}\cdot \L_1+\zeta_{\lambda} $ for some orders $\L_1$ and $\K_1$,
\item $\Z^{\lambda}\cdot \K_1+\zeta_{\lambda}\geq_\lambda \zeta_{\lambda} $ for some order $\K_1$,
\item $\zeta^*_{\lambda} +\Z^{\lambda}\cdot \K_1\geq_\lambda\zeta^*_{\lambda} $ for some order $\K_1$.
\end{enumerate}

The first three cases are handled immediately by $A_{\lambda}$. By symmetry it is clear that it is enough to show case 5. However, this is the same as showing that for all $\gamma<\lambda$,
$$\zeta_{\gamma}^* + \Z^\gamma\cdot(\omega +  \Z\cdot\zeta_{\lambda}^*)
\leq_{\gamma}\zeta_{\gamma}^* + \Z^{\gamma}\cdot
(\Z\cdot\zeta_{\lambda}^*+\Z^\lambda\cdot \K_1),$$
which follows from the $C_\gamma$ for $\gamma<\lambda$.

Therefore, for any countable ordinal $\alpha$, $A_\alpha$, $B_\alpha$ and $C_\alpha$ all hold.
\end{proof}

\end{document}
